\chapter{Administración de Mercados}
\label{ch:AdministracionMercados}

En el presente capítulo se ven aspectos relacionados con la estrategia mercadotécnica y se divide en tres partes: posicionamiento, promoción y precio.

El \emph{Posicionamiento} describe la identidad de la escuela. En la \emph{Promoción} se trazan las estrategias sugeridas para despertar el deseo entre la población de que los jóvenes se inscriban a la escuela. Por último, en la sección \emph{Precio} se mencionan algunas acciones de promoción y descuento específicas. 

\section{Posicionamiento}

\subsection{Identidad de la Escuela}

Hablar de la identidad de la escuela significa hablar de la persona del fundador. No es posible concebir la institución educativa salesiana sin considerar la figura de San Juan Bosco.

Por lo tanto, al proponer la imagen e identidad del bachillerato en la comunidad, el centro y punto de partida se centra en la polifacética y rica personalidad de un sacerdote que consagró su vida al servicio de los jóvenes, especialmente los más pobres.

\begin{figure}
	\centering
	\includegraphics[scale=0.5]{images/san-juan-bosco}
	\caption{San Juan Bosco, patrono de los jóvenes}
	\label{fig:don-bosco}
\end{figure}

Con motivo del 150 aniversario de la fundación de la Congregación de San Francisco de Sales (25 de abril de 2010) y del 200 aniversario del nacimiento del santo (1815-2015), las reliquias de Don Bosco visitan las comunidades de todo el mundo. En México, la gira inició el 4 de agosto en la Ciudad de México y concluyó el 11 de septiembre en Tijuana, B.C.N.; estuvo en Irapuato el 23 de agosto en una velada que atrajo a 5 mil personas.

En todas las poblaciones que estuvo la concurrencia fue nutrida, despertando expectación y curiosidad por conocer su vida y obra; se piensa aprovechar la sensibilización de la comunidad local para difundir información sobre el patrono de los jóvenes.

Algunos rasgos de su personalidad y su obra son:

\begin{itemize}
	\item Se le veía siempre rodeado de muchachos en un ambiente festivo, con música, risas y entusiasmo; conquistó el corazón de los jóvenes para educarlos.
	\item Instaló en sus escuelas diversos talleres (zapatería, sastrería, carpintería, imprenta) para preparar a los jóvenes para un trabajo digno.
	\item Confiaba en las persona y sabía ganarse su confianza, en una ocasión salió de paseo con los presos de la ciudad de Turín sin guardia alguno, al volver no escapó ninguno.
	\item Caminaba en la cuerda floja, realizaba malabares y trucos de magia, es patrono de los ilusionistas.
	\item Hospedaba huérfanos y les daba educación.
	\item Procuró mejorar las condiciones de trabajo de sus jóvenes y cuando no pudo convencer a sus patrones, él mismo les dio trabajo instalando talleres e imprenta.
\end{itemize}

Aplicado a la institución, significa considerar las notas señaladas en la sección \ref{sub:intro:NotasDistintivas} (página \pageref{sub:intro:NotasDistintivas}): \emph{sistema preventivo, acompañamiento, formación en valores, quitar barreras de discriminación, formación para el trabajo}

De este centro se desprenden de forma natural diversas características considerando los perfiles de clientes de la sección \ref{sec:Clientes} (página \pageref{sec:Clientes}) y en respuesta a la oferta de los competidores (cuadro \ref{tbl:CompetidoresDetalle}. p\'agina \pageref{tbl:CompetidoresDetalle}):

\begin{itemize}
	\item
	La preparatoria particular con el programa de becas más extenso, que busca continuamente cómo bajar sus precios para que nadie se quede sin estudiar.
	\item
	Uno de los colegios particulares de calidad con colegiaturas más bajas.
	\item
	Incorporada a la red de escuelas salesianas con todo el prestigio que implica.
	\item
	Atención y acompañamiento personalizados con base en el sistema preventivo.
	\item
	Uno de los pocos bachilleratos con carrera técnica, que capacita para el trabajo, para que el egresado pueda sostener sus estudios universitarios trabajando a tiempo parcial o si lo decide, incorporarse al medio laboral.
	\item
	La escuela tiene un fuerte compromiso en transmitir los valores cristianos y humanos formando <<Buenos Cristianos y Honestos Ciudadanos>> capaces de asumir compromisos con su familia y comunidad.
	\item
	El bachillerato tiene programas de vinculación empresarial con lo que permite la posibilidad de conseguir empleo a tiempo parcial o completo, conocer la industria y saber que la carrera técnica elegida corresponde con las necesidades laborales del mundo real.
\end{itemize}

\section{Promoción}

La promoción abarca cinco estrategias fundamentales, cada una de las cuales combina diversos instrumentos, se enfoca en determinados espacios y/o lugares, participan distintos agentes y persiguen objetivos precisos.

Los objetivos de la promoción son:

\begin{itemize}
	\item Interesar a los egresados de secundaria en ingresar al bachillerato tecnológico.
	\item Dar a conocer entre los vecinos, principalmente de los barrios populares, la nueva escuela como opción de superación.
	\item Dar a conocer al público en general la vida y obra de Don Bosco, su sistema educativo preventivo y la preparatoria que lo llevará a cabo en Irapuato.
	\item Llegar al público joven a través de Internet.
	\item Proporcionar informes y orientación.
\end{itemize}

Para lograrlo se tienen las siguientes cinco estrategias:

\begin{description}
	\item[Campaña de orientación vocacional en secundarias] visitas, conferencias, videos, carteles.
	\item[Campaña educativa en la calle] carteles en comercios, casas de la cultura, iglesias, mamparas en espacios públicos, etc. Volanteo en los barrios. Acceder a las comunidades de vecinos a través de la Iglesia en las misas, mediante grupos juveniles, aprovechando el apoyo de los párrocos, etc.
	\item[Promoción abierta] participación en programas de radio y televisión local, artículos y anuncios en el periódico, publicidad en espectaculares, etc.
	\item[Campaña por Internet] creación de sitio web, blog y cuentas en redes sociales (Facebook y Twitter). Elaboración de presentaciones electrónicas y colocación de videos en YouTube.
	\item[Centro telefónico] para proporcionar informes y orientación.
\end{description}

Las cinco estrategias se basan en la conformación de una red social integrada por miembros de la comunidad salesiana, ex-alumnos salesianos y voluntarios.

\section{Precio}

El precio base es de \$ 3,000.00. Sobre el precio existen diferentes estrategias:

\begin{itemize}
	\item Descuento por pago puntual
	\item Beca por aprovechamiento
	\item Beca por situación socioeconómica
	\item Beca por varios miembros de familia
	\item El porcentaje del monto asignado a becas se incrementa a medida que la matrícula crece
	\item Las becas se incrementan en la medida de los patrocinios que se obtengan
	\item Precio especial para hijos de trabajadores de empresas patrocinadoras
\end{itemize}
